% Source: https://artofproblemsolving.com/community/c3h1393998_mathcounts_state_target
% Stable identity: aops-topic-1393998-q01
Write $T_n=\dfrac{1}{2}+\dfrac{1}{3}+\cdots+\dfrac{1}{n}$; it increases with $n$, so we must find where it first exceeds $2$. We show $T_{10}<2<T_{11}$.

\textbf{Grouping.} Since $\dfrac{1}{2}+\dfrac{1}{3}+\dfrac{1}{6}=1$ and $\dfrac{1}{4}+\dfrac{1}{5}+\dfrac{1}{10}=\dfrac{5+4+2}{20}=\dfrac{11}{20}$,
\[
T_{10}=1+\dfrac{11}{20}+\left(\dfrac{1}{7}+\dfrac{1}{8}+\dfrac{1}{9}\right).
\]

\textbf{$T_{10}<2$.} Each of $\dfrac{1}{7},\dfrac{1}{8},\dfrac{1}{9}$ is at most $\dfrac{1}{7}$, so their sum is less than $\dfrac{3}{7}$, and $\dfrac{3}{7}<\dfrac{9}{20}$ because $60<63$. Hence $T_{10}<1+\dfrac{11}{20}+\dfrac{9}{20}=2$.

\textbf{$T_{11}>2$.} Now
\[
T_{11}=1+\dfrac{11}{20}+\left(\dfrac{1}{7}+\dfrac{1}{8}\right)+\left(\dfrac{1}{9}+\dfrac{1}{11}\right)
=1+\dfrac{11}{20}+\dfrac{15}{56}+\dfrac{20}{99}.
\]
Since $\dfrac{15}{56}>\dfrac{15}{60}=\dfrac{1}{4}$ and $\dfrac{20}{99}>\dfrac{20}{100}=\dfrac{1}{5}$, we get $T_{11}>1+\dfrac{11}{20}+\dfrac{1}{4}+\dfrac{1}{5}=1+\dfrac{11+5+4}{20}=2$.

Therefore the smallest such $n$ is $\boxed{11}$. (In fact $T_{10}\approx1.929$ and $T_{11}\approx2.020$.)

\emph{Remark.} The divergence of the infinite harmonic series is not needed here: the question is about a finite partial sum. Adding $1$ to both sides turns the condition into $1+\dfrac{1}{2}+\cdots+\dfrac{1}{n}>3$, a statement about the $n$th partial sum, not about $\displaystyle\sum_{n=1}^{\infty}\dfrac{1}{n}$.
